\chapter{Cerebral Aneurysm Clipping}
\section*{Overview}
Cerebral aneurysm clipping places a small clip across the neck of an intracranial aneurysm to prevent blood from entering it.
\section*{Common Instruments and Equipment}
The exact setup varies by technique and setting. Common items may include:
\begin{itemize}[nosep,leftmargin=*]
  \item craniotomy instruments and operating microscope
  \item temporary vascular clips and aneurysm clips
  \item micro-Doppler or intraoperative angiography when used
\end{itemize}

\section*{Why It Is Done}
It may treat a ruptured aneurysm or a selected unruptured aneurysm when open microsurgery is preferred to endovascular treatment.
\section*{How It Is Performed}
Through a craniotomy, the surgeon exposes the aneurysm and places a clip while preserving nearby arteries and nerves. Intraoperative imaging may confirm exclusion of the aneurysm.
\section*{Contraindications}
No single contraindication applies to every patient; the treating team assesses whether the procedure is suitable for the individual and the reason for treatment.

\section*{Precautions}
The care team reviews relevant health conditions, medicines, allergies, and procedure-specific risks, and explains any changes needed for the individual.

\section*{Risks and Recovery}
Stroke, bleeding, seizures, infection, brain swelling, cranial-nerve injury, and death are possible. Recovery depends heavily on whether the aneurysm ruptured before treatment.
\section*{References}
\begin{enumerate}\item National Institute of Neurological Disorders and Stroke. Cerebral Aneurysms. 2025.\item Current Medical Diagnosis and Treatment. McGraw-Hill; 2025.\end{enumerate}
\clearpage
